\chapter{Magnetostatics}

Earlier, we discussed electrostatics, where charges do not move, i.e.,
\begin{equation}
    \pdv{q}{t} = 0.
\end{equation}
Now we discuss \textbf{magnetostatics},
with the assumption of no electric forces, i.e.,
\begin{equation}
    \begin{aligned}
        Q &= 0, \\
        \vec{p} &= 0. \\
    \end{aligned}
\end{equation}

Magnetostatics examines the magnetic fields
produced by \textbf{steady currents}.
Unlike with electrostatics, there is no "magnetic charge",
or scalar source of magnetic fields.
More formally, this is given by \textbf{Gauss' law for magnetism}, which states
\begin{equation}
    \nabla \cdot \vec{B} = \vec{0}.
\end{equation}
This is sometimes called "the law of no magnetic monopoles",
i.e., there's no such thing as a north or south pole in isolation, only dipoles.

If we run two steady currents parallel to each other,
they will experience a attractive force.
If we run the currents antiparallel to each other,
they will experience a repulsive force.
